% CCSS 7.SP.A.1 — Populations and Samples (OK/TX: experimental design emphasis)
\section{Populations and Samples}

\begin{learningGoals}
\begin{itemize}
    \item Understand populations, samples, and why random sampling matters
    \item Design simple experiments and plan data collection
\end{itemize}
\end{learningGoals}

\begin{conceptBox}{Population, Sample, and Experimental Design}
\begin{itemize}[leftmargin=*]
    \item \textbf{Population}: The whole group you want to study.
    \item \textbf{Sample}: A smaller group chosen from the population.
    \item \textbf{Random sample}: Every member has an equal chance of being chosen. This avoids \textbf{bias}.
\end{itemize}
Before collecting data, plan your experiment:
\begin{enumerate}[leftmargin=*]
    \item \textbf{Ask a question} — What do you want to find out?
    \item \textbf{Identify the population} — Who or what are you studying?
    \item \textbf{Choose a sampling method} — How will you pick your sample?
    \item \textbf{Collect data} — Gather the information.
    \item \textbf{Analyze and conclude} — What does the data tell you?
\end{enumerate}
\end{conceptBox}

\begin{workedExample}{Designing a Survey}
A school wants to know the favorite sport of all $600$ students. Design a fair experiment.

\textbf{Question:} What sport is most popular among students?

\textbf{Population:} All $600$ students.

\textbf{Method:} Use a random number generator to pick $60$ student IDs. Survey those $60$.

\textbf{Why random?} Surveying only basketball players would bias the results. Random selection gives every student an equal chance.
\answerHighlight{A random sample of $60$ from $600$ avoids bias.}
\end{workedExample}

\tipBox{A sample must be random and representative to give valid results. Bigger samples are usually more reliable, but even a small random sample beats a large biased one.}

\begin{practiceBox}{Populations and Samples}
\resetProblems

\prob A teacher wants to know how much time students spend on homework. She surveys the first $10$ students who arrive. Is this a good sample? Explain.
\answerExplain{No — it is biased}{Students who arrive early may have different homework habits than others. A better method is to randomly select $10$ names from the class list.}

\prob Describe the population and sample: A factory tests $50$ light bulbs from a batch of $5{,}000$.
\answerExplain{Population: $5{,}000$ bulbs; sample: $50$ bulbs}{The factory uses the sample to estimate the quality of the whole batch without testing every single bulb.}

\prob Design an experiment to find out which school lunch is most popular. List the five steps.
\answerExplain{Question, population, method, collect, analyze}{1) Question: Which lunch is most popular? 2) Population: all students. 3) Method: randomly select $50$ students. 4) Collect: survey their lunch choice. 5) Analyze: tally results and find the most common choice.}

\prob Why is random sampling important?
\answerExplain{It avoids bias}{Random sampling gives every member an equal chance of being selected, so the sample is more likely to represent the whole population accurately.}

\prob A company surveys only its own employees about a new product. What is the problem?
\answerExplain{Biased sample}{Employees may have opinions influenced by loyalty or job concerns. The population of interest is all potential customers, not just employees.}

\end{practiceBox}
